% problem_id: S001-lemma-completion-countably-indexed
% source_id: S001 (Stacks Project)
% source_file: formal-spaces.tex
% statement_locator: formal-spaces.tex lines 3094--3108
% proof_locator: formal-spaces.tex lines 3110--3229
% env: lemma
% id_note: label 在本来源内唯一
% pairing: tight (verified)
% status: complete (statement + author proof verbatim + reason + locators)
%
\begin{lemma}
\label{lemma-completion-countably-indexed}
\begin{reference}
Email by Ofer Gabber of September 11, 2014.
\end{reference}
Let $S$ be a scheme. Let $X = \Spec(A)$ be an affine scheme over $S$.
Let $T \subset X$ be a closed subscheme.
\begin{enumerate}
\item If the formal completion $X_{/T}$ is countably indexed
and there exist countably many $f_1, f_2, f_3, \ldots \in A$ such that
$T = V(f_1, f_2, f_3, \ldots)$, then $X_{/T}$ is adic*.
\item The conclusion of (1) is wrong if we omit the assumption that
$T$ can be cut out by countably many functions in $X$.
\end{enumerate}
\end{lemma}

\begin{proof}
The assumption that $X_{/T}$ is countably indexed means that there exists a
sequence of ideals
$$
A \supset J_1 \supset J_2 \supset J_3 \supset \ldots
$$
with $V(J_n) = T$ such that every ideal $J \subset A$ with $V(J) = T$
there exists an $n$ such that $J \supset J_n$.

\medskip\noindent
To construct an example for (2) let $\omega_1$ be the first uncountable
ordinal. Let $k$ be a field and let
$A$ be the $k$-algebra generated by $x_\alpha$, $\alpha \in \omega_1$
and $y_{\alpha \beta}$ with $\alpha \in \beta \in \omega_1$
subject to the relations $x_\alpha = y_{\alpha \beta} x_\beta$.
Let $T = V(x_\alpha)$. Let $J_n = (x_\alpha^n)$.
If $J \subset A$ is an ideal such that
$V(J) = T$, then $x_\alpha^{n_\alpha} \in J$ for some $n_\alpha \geq 1$.
One of the sets $\{\alpha \mid n_\alpha = n\}$ must be unbounded in
$\omega_1$. Then the relations imply that $J_n \subset J$.

\medskip\noindent
To see that (2) holds it now suffices to show that $A^\wedge = \lim A/J_n$
is not a ring complete with respect to a finitely generated ideal.
For $\gamma \in \omega_1$ let $A_\gamma$ be the quotient of $A$
by the ideal generated by $x_\alpha$, $\alpha \in \gamma$ and
$y_{\alpha \beta}$, $\alpha \in \gamma$. As $A/J_1$ is reduced,
every topologically nilpotent element $f$ of $\lim A/J_n$ is in
$J_1^\wedge = \lim J_1/J_n$. This means $f$ is an infinite series
involving only a countable number of generators. Hence $f$ dies in
$A_\gamma^\wedge = \lim A_\gamma/J_nA_\gamma$ for some $\gamma$.
Note that $A^\wedge \to A_\gamma^\wedge$ is continuous and open by
Lemma \ref{lemma-ses}.
If the topology on $A^\wedge$ was $I$-adic for some finitely generated ideal
$I \subset A^\wedge$, then $I$ would go to zero in some
$A_\gamma^\wedge$. This would mean that $A_\gamma^\wedge$ is discrete,
which is not the case as there is a surjective continuous and open
(by Lemma \ref{lemma-ses}) map
$A_\gamma^\wedge \to k[[t]]$ given by
$x_\alpha \mapsto t$, $y_{\alpha \beta} \mapsto 1$ for
$\gamma = \alpha$ or $\gamma \in \alpha$.

\medskip\noindent
Before we prove (1) we first prove the following: If $I \subset A^\wedge$ is
a finitely generated ideal whose closure $\bar I$ is open, then $I = \bar I$.
Since $V(J_n^2) = T$ there exists an $m$ such that $J_n^2 \supset J_m$.
Thus, we may assume that $J_n^2 \supset J_{n + 1}$ for all $n$ by passing
to a subsequence. Set $J_n^\wedge = \lim_{k \geq n} J_n/J_k \subset A^\wedge$.
Since the closure $\bar I = \bigcap (I + J_n^\wedge)$
(Lemma \ref{lemma-closed}) is open we see that there exists an $m$ such that
$I + J_n^\wedge \supset J_m^\wedge$ for all $n \geq m$. Fix such an $m$.
We have
$$
J_{n - 1}^\wedge I + J_{n + 1}^\wedge \supset
J_{n - 1}^\wedge (I + J_{n + 1}^\wedge) \supset
J_{n - 1}^\wedge J_m^\wedge
$$
for all $n \geq m + 1$. Namely, the first inclusion is trivial and the
second was shown above. Because $J_{n - 1}J_m \supset J_{n - 1}^2 \supset J_n$
these inclusions show that the image of $J_n$ in $A^\wedge$
is contained in the ideal $J_{n - 1}^\wedge I + J_{n + 1}^\wedge$.
Because this ideal is open we conclude that
$$
J_{n - 1}^\wedge I + J_{n + 1}^\wedge \supset J_n^\wedge.
$$
Say $I = (g_1, \ldots, g_t)$. Pick $f \in J_{m + 1}^\wedge$.
Using the last displayed inclusion, valid for all $n \geq m + 1$,
we can write by induction on $c \geq 0$
$$
f = \sum f_{i, c} g_i \mod J_{m + 1+ c}^\wedge
$$
with $f_{i, c} \in J_m^\wedge$ and
$f_{i, c} \equiv f_{i, c - 1} \bmod J_{m + c}^\wedge$.
It follows that $IJ_m^\wedge \supset J_{m + 1}^\wedge$.
Combined with $I + J_{m + 1}^\wedge \supset J_m^\wedge$
we conclude that $I$ is open.

\medskip\noindent
Proof of (1). Assume $T = V(f_1, f_2, f_3, \ldots)$.
Let $I_m \subset A^\wedge$ be the ideal generated by $f_1, \ldots, f_m$.
We distinguish two cases.

\medskip\noindent
Case I: For some $m$ the closure of $I_m$ is open.
Then $I_m$ is open by the result of the previous paragraph.
For any $n$ we have $(J_n)^2 \supset J_{n+1}$ by design, so
the closure of $(J_n^\wedge)^2$ contains $J_{n+1}^\wedge$
and thus is open.  Taking $n$ large, it follows that the closure
of the product of any two open ideals in $A^\wedge$ is open.
Let us prove $I_m^k$ is open for $k \ge 1$ by induction on $k$.
The case $k = 1$ is our hypothesis on $m$ in Case I.
For $k > 1$, suppose $I_m^{k - 1}$ is open. Then
$I_m^k = I_m^{k - 1} \cdot I_m$ is the product of two open ideals
and hence has open closure. But then since $I_m^k$
is finitely generated it follows that $I_m^k$
is open by the previous paragraph (applied to $I = I_m^k$),
so we can continue the induction on $k$.
As each element of $I_m$ is topologically nilpotent, we conclude
that $I_m$ is an ideal of definition which proves that $A^\wedge$
is adic with a finitely generated ideal of definition, i.e.,
$X_{/T}$ is adic*.

\medskip\noindent
Case II. For all $m$ the closure $\bar I_m$ of $I_m$ is not open.
Then the topology on $A^\wedge/\bar I_m$ is not discrete. This means
we can pick $\phi(m) \geq m$ such that
$$
\Im(J_{\phi(m)} \to A/(f_1, \ldots, f_m)) \not =
\Im(J_{\phi(m) + 1} \to A/(f_1, \ldots, f_m))
$$
To see this we have used that
$A^\wedge/(\bar I_m + J_n^\wedge) = A/((f_1, \ldots, f_m) + J_n)$.
Choose exponents $e_i > 0$ such that $f_i^{e_i} \in J_{\phi(m) + 1}$
for $0 < m < i$. Let $J = (f_1^{e_1}, f_2^{e_2}, f_3^{e_3}, \ldots)$.
Then $V(J) = T$. We claim that $J \not \supset J_n$ for all $n$
which is a contradiction proving Case II does not occur.
Namely, the image of $J$ in $A/(f_1, \ldots, f_m)$ is contained
in the image of $J_{\phi(m) + 1}$ which is properly contained in the
image of $J_m$.
\end{proof}
